\chapter{Concepts}
\label{ch:concepts}

This chapter is the model of the kernel: what a task is, when it runs, how time is
counted, what an interrupt may do, and how mutual exclusion works. Everything in the
later chapters assumes it.

\section{Tasks and priorities}

A task is a C function that never returns, running on a stack the application provides,
with a priority that is also its identifier. There are \code{RTOS\_CFG\_NUM\_TASKS}
application tasks, numbered 1 to \code{RTOS\_CFG\_NUM\_TASKS} (at most 31), and the idle
task 0, which is the context that called \code{rtos\_start}. A higher number is a higher
priority. Because the id is the priority, every kernel object that records waiting tasks
does so as a bit mask (\code{rtos\_taskmask\_t}), and the highest waiter is the highest
set bit.

The task control block (TCB) holds the saved context, the stack, the pending event flags,
the wait state, the timeout and the statistics; it is 42 bytes on the AVR and 72 on the
ESP32 in the default configuration (figure~\ref{fig:ram}). The kernel adds nothing per
task at run time.

\section{Scheduling}
\label{sec:sched}

The scheduler keeps one \emph{ready mask} with a bit per task that is ready to run. The
running task's bit stays set. The next task to run is the highest set bit: a lookup
table on the AVR, one instruction on the Xtensa. A task that blocks clears its bit and
the kernel switches to the highest remaining one; a task that is woken sets its bit, and
if it is higher than the running task, the kernel switches to it at once (from task
context) or at the end of the interrupt (from an ISR body).

The rule an application can rely on:

\begin{quote}
At every instant the running task is the ready task with the greatest key
$(\mathit{effPrio}, \mathit{boosted}, \mathit{prio})$, compared lexicographically. A
ready task preempts the running one only if its key is strictly greater.
\end{quote}

Without mutexes the effective priority equals the priority, no task is boosted, and the
rule reduces to \emph{the highest priority ready task runs}. The two other components
matter only while a task holds a mutex whose ceiling is above its own priority
(section~\ref{sec:mutex}); because priorities are unique, the order is total and no two
tasks ever tie. There is no time slicing: a task runs until it blocks, is preempted by a
higher task, or returns.

The idle task runs when no other task is ready. It executes the idle hook, the watchdog
and the stack monitor when they are configured, and then, with
\code{RTOS\_CFG\_IDLE\_SLEEP}, puts the CPU to sleep until the next interrupt. With
\code{RTOS\_CFG\_TICKLESS} on the ESP32 it also stops the tick and sleeps until the next
deadline. The idle task must never block; a blocking call from it is checked.

\section{Time}
\label{sec:time}

The tick is a periodic interrupt at \code{RTOS\_CFG\_TICK\_HZ}, 1\,kHz in every shipped
configuration. It increments the tick counter, wakes the tasks whose timeout has
expired (they are kept in a list sorted by wake-up time, so the tick looks at the head
only) and runs the software timers that are due. \code{rtos\_getTick} returns the
counter; it is 16 or 32 bits wide (\code{RTOS\_CFG\_TICK\_BITS}) and wraps, so tick
values are always compared as differences, never directly. The kernel does the same, and
that is what limits every finite timeout to half the counter's range,
\code{RTOS\_MAX\_TIMEOUT}: 32767 ticks with 16 bit ticks, $2^{31}-1$ with 32. A longer
timeout is checked; \code{rtos\_delay} is the one exception and splits a long delay into
pieces.

Two things in the semantics of time are easy to get wrong:

\begin{itemize}
\item \textbf{A delay of $n$ ticks lasts between $n$ and $n+1$ tick periods}, because
  part of the current period is already over when the call is made. \code{rtos\_delay(0)}
  waits for the next tick; it is not a yield (there is nothing to yield to: a task of the
  same priority does not exist).
\item \textbf{A periodic task does not drift.} \code{rtos\_waitPeriodic(period)} waits
  for the next boundary, and each boundary is the previous one plus the period, not the
  time of the call plus the period. If the boundary has passed, the call returns at once,
  counts an overrun and starts the next period from now. A task whose first iteration
  takes longer than a period therefore reports one overrun at start-up, harmlessly.
\end{itemize}

\code{rtos\_getTimeUs} is a finer clock, read from a free running hardware counter
(4\,$\mu$s steps on the AVR, the cycle counter on the ESP32); it wraps after 71.6 minutes
and is for measuring intervals.

\section{Interrupts and calling contexts}
\label{sec:contexts}

An interrupt handler that talks to the kernel is written with the \code{RTOS\_ISR}
macro:

\begin{lstlisting}
RTOS_ISR(BOARD_APP_TIMER_VECT)          /* a timer compare interrupt */
{
    rtos_semGiveFromIsr(&_sem);         /* wakes the consumer, no switch yet */
}
\end{lstlisting}

On the AVR the macro produces a naked handler that saves the caller clobbered registers,
runs the body and ends in the kernel's \emph{interrupt tail}: if the body made a higher
task ready, the tail saves the rest of the context and switches, otherwise it returns to
the interrupted task. On the ESP32 the macro defines the body as a function the port's
exception dispatcher calls for that interrupt, and the tail is inlined into the
dispatcher. In both cases the rule is the same: an ISR body never switches; \emph{the
switch happens once, at the end of the ISR}. That is what makes several wake-ups in one
interrupt cost one switch, and it is why there are separate \code{FromIsr} functions:
they do the same work as the task variants but defer the switch.

Interrupts do not nest on either port; a second interrupt waits and runs on the same
spot. The longest interrupt-off window is therefore a context switch (about 27\,$\mu$s on
the Mega, 1.7\,$\mu$s on the ESP32), and that is the worst case interrupt latency the
benchmark in chapter~\ref{ch:perf} reports.

\begin{table}[htbp]
\centering\small
\begin{tabularx}{\linewidth}{@{}lX@{}}
\toprule
Context & May call \\
\midrule
Task & everything without a \code{FromIsr} suffix; a call may block or switch at once \\
ISR body, timer callback in the tick, trace hook & the \code{FromIsr} variants,
  \code{rtos\_timerStart}, \code{rtos\_timerStop}, \code{rtos\_getTick},
  \code{rtos\_getTimeUs}, the readers (\code{rtos\_queueCount}, \code{rtos\_mboxIsFull},
  \code{rtos\_poolFreeCount}, \code{rtos\_evgroupGet}, \code{rtos\_timerIsActive},
  \code{rtos\_timerRemaining}) and \code{rtos\_panic}; nothing blocks or switches \\
Idle task, idle hook & as a task, but nothing that blocks: waits only with
  \code{RTOS\_NO\_WAIT}; a blocking call is checked \\
Timer callback in the timer task & everything a task may call \\
Panic hook & nothing of the kernel; interrupts are disabled \\
\bottomrule
\end{tabularx}
\caption{Calling contexts.}
\label{tab:contexts}
\end{table}

\paragraph{Critical sections.} The kernel's own critical sections are short interrupt-off
windows (\code{rtos\_port\_enterCritical} returns the previous state,
\code{rtos\_port\_exitCritical} restores it, so they nest). An application may use the
same pair for a few instructions of its own; it costs three cycles on the AVR and six on
the ESP32. Everything longer belongs behind a mutex or inside a single task.

\section{Blocking and hand-over}

Every blocking call takes a timeout: \code{RTOS\_NO\_WAIT} polls, \code{RTOS\_FOREVER}
waits for good, anything else is ticks up to \code{RTOS\_MAX\_TIMEOUT}. A blocked task
records what it waits on (\code{waitKind}, visible in \code{rtos\_getTaskInfo}) and, if
it has a timeout, sits in the timed list.

When a resource becomes available and a task waits for it, the kernel \emph{hands it
over} directly: a semaphore token goes to the highest waiter without touching the count,
a queue item is copied straight into the receiver's buffer, a freed pool block goes to
the waiting allocator, a mutex is passed to the highest waiter who becomes the owner at
once. The resource never passes through the shared state, so a lower task cannot sneak
in between the release and the wake-up, and the waiter's return value is decided at the
moment of the release. The woken task runs before the releasing call returns if it is
higher.

\section{Mutual exclusion}
\label{sec:mutex}

\subsection{The problem}

A shared resource guarded by a plain lock produces \emph{priority inversion}: a high task
$H$ blocks on a lock held by a low task $L$, and a medium task $M$ that needs neither
preempts $L$ and keeps $H$ waiting for as long as it likes. The classic remedies are
priority inheritance, where $L$ is raised to $H$'s priority once $H$ blocks, and the
priority ceiling protocols, where the lock itself carries a priority~\cite{shaetal}.

\subsection{The immediate priority ceiling protocol}

nornCore implements the immediate variant, also known as the stack resource policy or
highest locker protocol~\cite{baker,burnswellings}. Every mutex has a \emph{ceiling}: the
highest priority of any task that will ever lock it. Locking the mutex raises the
caller's \emph{effective priority} to the ceiling for as long as it holds the mutex;
unlocking restores the previous one. The consequences, all of which the kernel relies
on:

\begin{itemize}
\item \textbf{Nobody blocks on a mutex.} While $L$ holds the mutex it runs at the
  ceiling, which is at least $H$'s priority, so $H$ cannot run and cannot find the mutex
  locked. The lock operation is a priority change and a few stores: 218 cycles for a
  lock and unlock pair on the Mega, 156 on the ESP32, with no scheduler involvement.
\item \textbf{Inversion is bounded by one critical section.} $H$ can be delayed by at
  most one lower task's critical section, and only one that was already running when $H$
  became ready.
\item \textbf{Deadlock is impossible} with correct ceilings: a task that holds one mutex
  runs at a priority at which no other user of any mutex it will take next can hold it.
\item \textbf{A medium task $M$ with priority below the ceiling is held off} while $L$
  holds the mutex, even if $M$ never uses it. That is the price of the protocol and the
  reason ceilings should be set exactly, not generously.
\end{itemize}

\begin{figure}[htbp]
\centering
\begin{tikzpicture}[x=0.9cm,y=0.9cm,font=\small]
  % axes
  \draw[-{Latex}] (0,0) -- (12.5,0) node[below left] {time};
  \foreach \y/\n in {1/L (1),2/M (2),3/H (3)} {
    \draw[gray!60] (0,\y) -- (12,\y);
    \node[left] at (0,\y) {\n};
  }
  % L runs, locks at t=1 (raised to 3), unlocks at t=5
  \fill[nornblue!30] (0,0.6) rectangle (1,1.4);
  \fill[nornblue!60] (1,2.6) rectangle (5,3.4);
  \draw[dashed] (1,1.4) -- (1,2.6);
  \draw[dashed] (5,2.6) -- (5,1.4);
  \node[above,font=\footnotesize] at (3,3.4) {L holds the mutex at the ceiling 3};
  % H ready at t=2, runs at t=5
  \draw[-{Latex},nornred,thick] (2,4.2) -- (2,3.5) node[midway,right,font=\footnotesize] {H ready};
  \fill[nornred!50] (5,2.6) rectangle (7.5,3.4);
  \node[above,font=\footnotesize] at (6.25,3.4) {H};
  % M ready at t=3, runs after H
  \draw[-{Latex},black!70,thick] (3,4.2) -- (3,3.5) node[midway,right,font=\footnotesize] {M ready};
  \fill[gray!40] (7.5,1.6) rectangle (9.5,2.4);
  \node[above,font=\footnotesize] at (8.5,2.4) {M};
  \fill[nornblue!30] (9.5,0.6) rectangle (11,1.4);
  \node[above,font=\footnotesize] at (10.25,1.4) {L};
  \foreach \x in {0,1,...,12} \draw (\x,-0.1) -- (\x,0.1);
\end{tikzpicture}
\caption{Immediate priority ceiling. L (priority 1) locks a mutex with ceiling 3 at $t=1$
and runs at 3 until the unlock at $t=5$. H (3) and M (2) become ready meanwhile and
wait; H is delayed by one critical section of L and by nothing else, and M, which never
uses the mutex, is held off as long as L holds it. Nobody ever blocks on the mutex
itself.}
\label{fig:ceiling}
\end{figure}

\subsection{When the ceiling is wrong, or the owner blocks}

Two things break the ideal picture. A ceiling that is \emph{too low} (a task above the
ceiling locks the mutex) is a configuration error; the lock is asserted against it, and
without assertions the task finds the mutex held. And an owner may \emph{block while
holding} the mutex (allowed, discouraged): it then leaves the CPU at the ceiling, and a
task above its base priority but at or below the ceiling can run and try the lock.

For both cases the kernel has a fallback: the caller waits on the mutex, without a
timeout, and the unlock hands the mutex to the highest waiter, who becomes the owner at
the ceiling at once. A wait that would close a cycle of tasks waiting for each other's
mutexes is detected and is \code{RTOS\_PANIC\_DEADLOCK}, naming both tasks. The detector
follows chains of mutex waits only, which is why a task must not return while holding a
mutex (section~\ref{sec:lifecycle}): an owner that is gone would make every later lock
wait for ever without a cycle to detect.

\subsection{The order at equal effective priority}
\label{sec:tie}

While a task holds a mutex at the ceiling $c$, the task whose base priority is $c$ may
become ready. Both have effective priority $c$. The scheduling key from
section~\ref{sec:sched} settles it: the \emph{boosted} owner ranks above the task whose
base priority is the ceiling, so the owner finishes its critical section first and the
higher task does not preempt it. This is the choice that preserves the protocol's
promise that nobody blocks: were the task at the ceiling to preempt, it would find the
mutex locked. Two owners at one ceiling, which is possible only if one of them blocked
while holding, run in base priority order. The test \code{test/test\_tie.c} pins every
one of these cases down.

\subsection{Rules for using mutexes}

\begin{enumerate}
\item Set the ceiling to the highest priority of any task that locks the mutex, exactly.
\item Keep the critical section short; nothing below the ceiling runs meanwhile.
\item Do not block while holding a mutex. It is allowed, and the kernel copes, but it
  reintroduces waiting and it is the only way to get two tasks at one ceiling.
\item Lock in a fixed order and unlock in the reverse order (asserted). A mutex is not
  recursive (asserted).
\item Never return from a task function while holding a mutex (checked for every mutex
  whose ceiling is above the task's priority).
\item Do not lock from the idle task (asserted) or from an ISR body: an ISR body cannot
  wait, and a critical section is the right tool for a few instructions of shared state
  with an interrupt.
\end{enumerate}

\section{The task life cycle}
\label{sec:lifecycle}

A task slot is in one of five states, derived from the TCB rather than stored in it (no
state byte, no store on the hot path), with suspension as a modifier on top of READY and
BLOCKED:

\begin{table}[htbp]
\centering\small
\begin{tabularx}{\linewidth}{@{}lXll@{}}
\toprule
State & Predicate & Entered by & Left by \\
\midrule
UNUSED     & no task function; the TCB is in its initial state & \code{rtos\_init}, a future reap & \code{rtos\_createTask} \\
READY      & bit in the ready mask, not the running task & create, wake-up, preemption & pick, block, suspend \\
RUNNING    & the current task (its bit stays in the ready mask) & pick & block, preemption, return \\
BLOCKED    & bit not in the ready mask; \code{waitKind} says on what & a blocking call & wake-up, timeout \\
TERMINATED & \code{waitKind == RTOS\_WAIT\_TERMINATED} & the task function returning & a future reap \\
\bottomrule
\end{tabularx}
\caption{The task states. \code{rtos\_getTaskInfo} reports them as \code{RTOS\_TASK\_*},
with SUSPENDED for a suspended READY or BLOCKED task.}
\label{tab:states}
\end{table}

\begin{figure}[htbp]
\centering
\begin{tikzpicture}[font=\small, node distance=2.2cm and 2.6cm,
  state/.style={draw, rounded corners, minimum width=2.3cm, minimum height=0.8cm, fill=codebg},
  >={Latex[length=2mm]}]
  \node[state] (unused) {UNUSED};
  \node[state, right=of unused] (ready) {READY};
  \node[state, right=of ready] (running) {RUNNING};
  \node[state, below=of running] (blocked) {BLOCKED};
  \node[state, below=of ready] (term) {TERMINATED};
  \draw[->] (unused) -- node[above] {createTask} (ready);
  \draw[->] (ready.north east) to[bend left=20] node[above] {pick} (running.north west);
  \draw[->] (running.south west) to[bend left=20] node[below] {preempted} (ready.south east);
  \draw[->] (running) -- node[right] {blocking call} (blocked);
  \draw[->] (blocked) -- node[below, sloped] {wake-up, timeout} (ready);
  \draw[->] (running) -- node[above, sloped, pos=0.6] {function returns} (term);
  \draw[->, dashed] (term) -- node[left, text width=1.9cm, align=center] {reap\\(future)} (unused);
\end{tikzpicture}
\caption{The life cycle of a task slot. Suspension (\code{rtos\_suspend}) is a modifier
on READY and BLOCKED, not a state: the task keeps its wait and its pending wake-ups and
continues at \code{rtos\_resume}.}
\label{fig:states}
\end{figure}

A task function \emph{should} never return. If it does, the task is terminated: its slot
and stack stay allocated, it is never scheduled again, and it is deaf to events,
notifications, suspends and resumes aimed at it (they accumulate or are ignored
harmlessly). Returning while holding a mutex is checked, for the reason given above. The
kernel does not yet reuse a terminated slot at run time; the rules under which a future
dynamic task API will do so are fixed (appendix~\ref{app:roadmap}), and the test
\code{test/test\_lifecycle.c} already verifies them against the kernel.

\section{Stacks}
\label{sec:stacks}

Each task runs on the stack given to \code{rtos\_createTask}. What a task needs is its
own frames plus what the kernel and the port leave under them when the task blocks or is
interrupted. An interrupt lands on the stack of the running task, so every task must
have room for the interrupt frame and the kernel's work on top of its own deepest call
chain. The port gives three sizes (table~\ref{tab:stacksizes}), and chapter~\ref{ch:perf}
has the measurements they are derived from.

\begin{table}[htbp]
\centering\small
\begin{tabularx}{\linewidth}{@{}lrrX@{}}
\toprule
Name & ESP32 & AVR & Meant for \\
\midrule
\code{RTOS\_STACK\_MIN}     & 512 B  & 96 B  & the technical minimum the kernel accepts (checked): very simple tasks whose use has been measured \\
\code{RTOS\_STACK\_DEFAULT} & 1024 B & 192 B & an ordinary task; start here \\
\code{RTOS\_STACK\_LARGE}   & 2048 B & 384 B & a recommended start for tasks that print, nest calls or use floating point; not an upper bound \\
\bottomrule
\end{tabularx}
\caption{The stack sizes the port names. Size a task from the default, run it, read the
high water mark with \code{rtos\_getStackFree}, and keep a margin above
\code{RTOS\_PORT\_STACK\_WARN} (96 B on the ESP32, 24 B on the AVR).}
\label{tab:stacksizes}
\end{table}

With \code{RTOS\_CFG\_STACK\_CHECK} the stack is filled with a pattern at creation and
the high water mark can be read; with \code{RTOS\_CFG\_STACK\_WARN} the idle task
watches one task per loop and calls \code{rtos\_stackWarnHook} once for a task whose
margin is gone; with \code{RTOS\_CFG\_STACK\_CHECK 2} a canary at the bottom turns an
actual overflow into \code{RTOS\_PANIC\_STACK\_OVERFLOW}, checked before every switch and
at every tick.
